\documentclass[11pt,letterpaper,openany]{book}
\title{GSD: The Beginner's Guide}
\author{Tibsfox}
\date{}

% Tell XeLaTeX where to find shared inputs
\makeatletter
\def\input@path{{../shared/}}
\makeatother

% ============================================================================
%  gsd-guides-preamble.tex
%  Shared LaTeX preamble for the GSD User Guides Suite
%  Compile with: xelatex (3-pass), NOT pdflatex
%  License: CC BY-SA 4.0 — see \cclicense macro
%
%  Load order note: this file loads xcolor + tcolorbox BEFORE sourcing
%  gsd-color-scheme.tex, and then loads hyperref AFTER color definitions
%  are live. Documents only need to % ============================================================================
%  gsd-guides-preamble.tex
%  Shared LaTeX preamble for the GSD User Guides Suite
%  Compile with: xelatex (3-pass), NOT pdflatex
%  License: CC BY-SA 4.0 — see \cclicense macro
%
%  Load order note: this file loads xcolor + tcolorbox BEFORE sourcing
%  gsd-color-scheme.tex, and then loads hyperref AFTER color definitions
%  are live. Documents only need to % ============================================================================
%  gsd-guides-preamble.tex
%  Shared LaTeX preamble for the GSD User Guides Suite
%  Compile with: xelatex (3-pass), NOT pdflatex
%  License: CC BY-SA 4.0 — see \cclicense macro
%
%  Load order note: this file loads xcolor + tcolorbox BEFORE sourcing
%  gsd-color-scheme.tex, and then loads hyperref AFTER color definitions
%  are live. Documents only need to \input{gsd-guides-preamble} in their
%  preamble (this file pulls in the color scheme transitively).
% ============================================================================

% ---------- Fonts ----------
\usepackage{fontspec}
\setmainfont{DejaVu Serif}[
  Ligatures=TeX,
  Scale=0.95,
  BoldFont={DejaVu Serif Bold},
  ItalicFont={DejaVu Serif Italic},
  BoldItalicFont={DejaVu Serif Bold Italic}
]
\newfontfamily\headingfont{DejaVu Sans}[Scale=0.95]
\newfontfamily\monofont{DejaVu Sans Mono}[Scale=0.88]
\setmonofont{DejaVu Sans Mono}[Scale=0.88]

% ---------- Page geometry ----------
\usepackage[
  letterpaper,
  top=1in,
  bottom=1in,
  left=1.25in,
  right=1.25in,
  headheight=14pt
]{geometry}

% ---------- Color + callout engine (must load BEFORE color-scheme input) ----
\usepackage{xcolor}
\usepackage[most]{tcolorbox}

% ---------- Typesetting packages ----------
\usepackage{tabularx}
\usepackage{longtable}
\usepackage{booktabs}
\usepackage{colortbl}
\usepackage{lastpage}
\usepackage{enumitem}
\usepackage{listings}
\usepackage{tikz}
\usepackage{fancyhdr}
\usepackage{titlesec}
\usepackage{parskip}
\usepackage{microtype}

% ---------- Pull in the color scheme (defines gsdnavy, gsdaccent, callouts) -
\input{gsd-color-scheme.tex}

% ---------- Hyperref (needs colors already defined) ----------
\usepackage[
  colorlinks=true,
  linkcolor=gsdnavy,
  urlcolor=gsdaccent,
  citecolor=gsdnavy,
  bookmarksnumbered=true
]{hyperref}

% ---------- Header / footer ----------
\pagestyle{fancy}
\fancyhf{}
\fancyhead[L]{\small\headingfont\@title}
\fancyhead[R]{\small\headingfont\leftmark}
\fancyfoot[L]{\small\headingfont GSD v1.35.0 — CC BY-SA 4.0}
\fancyfoot[R]{\small\headingfont Page \thepage\ of \pageref*{LastPage}}
\renewcommand{\headrulewidth}{0.4pt}
\renewcommand{\footrulewidth}{0.4pt}

\fancypagestyle{plain}{%
  \fancyhf{}
  \fancyfoot[L]{\small\headingfont GSD v1.35.0 — CC BY-SA 4.0}
  \fancyfoot[R]{\small\headingfont Page \thepage\ of \pageref*{LastPage}}
  \renewcommand{\headrulewidth}{0pt}
  \renewcommand{\footrulewidth}{0.4pt}
}

% ---------- Chapter / section styling ----------
\titleformat{\chapter}[display]
  {\headingfont\huge\bfseries\color{gsdnavy}}
  {\chaptertitlename\ \thechapter}{20pt}{\Huge}
\titleformat{\section}
  {\headingfont\Large\bfseries\color{gsdnavy}}
  {\thesection}{1em}{}
\titleformat{\subsection}
  {\headingfont\large\bfseries\color{gsdnavy}}
  {\thesubsection}{1em}{}

% ---------- Listings (code blocks) ----------
\lstdefinestyle{gsdcode}{
  basicstyle=\ttfamily\small,
  backgroundcolor=\color{gsdcodebg},
  frame=leftline,
  framerule=2pt,
  rulecolor=\color{gsdaccent},
  xleftmargin=1em,
  xrightmargin=0.5em,
  aboveskip=0.8em,
  belowskip=0.8em,
  breaklines=true,
  breakatwhitespace=true,
  showstringspaces=false,
  columns=flexible,
  keepspaces=true
}
\lstset{style=gsdcode}

% ---------- Semantic macros ----------
\newcommand{\cclicense}{%
  \textit{Released under CC BY-SA 4.0.\\
  Authored by Tibsfox. GSD itself is MIT-licensed — see gsd-build/get-shit-done.}%
}

\newcommand{\chapterheading}[1]{{\headingfont\LARGE\bfseries\color{gsdnavy}#1\par\medskip}}
\newcommand{\sectionheading}[1]{{\headingfont\Large\bfseries\color{gsdnavy}#1\par\smallskip}}
 in their
%  preamble (this file pulls in the color scheme transitively).
% ============================================================================

% ---------- Fonts ----------
\usepackage{fontspec}
\setmainfont{DejaVu Serif}[
  Ligatures=TeX,
  Scale=0.95,
  BoldFont={DejaVu Serif Bold},
  ItalicFont={DejaVu Serif Italic},
  BoldItalicFont={DejaVu Serif Bold Italic}
]
\newfontfamily\headingfont{DejaVu Sans}[Scale=0.95]
\newfontfamily\monofont{DejaVu Sans Mono}[Scale=0.88]
\setmonofont{DejaVu Sans Mono}[Scale=0.88]

% ---------- Page geometry ----------
\usepackage[
  letterpaper,
  top=1in,
  bottom=1in,
  left=1.25in,
  right=1.25in,
  headheight=14pt
]{geometry}

% ---------- Color + callout engine (must load BEFORE color-scheme input) ----
\usepackage{xcolor}
\usepackage[most]{tcolorbox}

% ---------- Typesetting packages ----------
\usepackage{tabularx}
\usepackage{longtable}
\usepackage{booktabs}
\usepackage{colortbl}
\usepackage{lastpage}
\usepackage{enumitem}
\usepackage{listings}
\usepackage{tikz}
\usepackage{fancyhdr}
\usepackage{titlesec}
\usepackage{parskip}
\usepackage{microtype}

% ---------- Pull in the color scheme (defines gsdnavy, gsdaccent, callouts) -
% ============================================================================
%  gsd-color-scheme.tex
%  Color palette + tcolorbox environment definitions for all GSD User Guides
% ============================================================================

\definecolor{gsdnavy}{HTML}{1B2A4A}
\definecolor{gsdaccent}{HTML}{CB3837}
\definecolor{gsdcodebg}{HTML}{F5F5F5}
\definecolor{gsdcodeborder}{HTML}{CCCCCC}

\definecolor{gsdtipbg}{HTML}{E8F4FD}
\definecolor{gsdtipbar}{HTML}{2B6CB0}

\definecolor{gsdwarnbg}{HTML}{FFF3CD}
\definecolor{gsdwarnbar}{HTML}{B78D00}

\definecolor{gsddangerbg}{HTML}{F8D7DA}
\definecolor{gsddangerbar}{HTML}{B02A37}

\definecolor{gsdsuccessbg}{HTML}{D4EDDA}
\definecolor{gsdsuccessbar}{HTML}{198754}

\definecolor{gsdnotebg}{HTML}{F0F0F5}
\definecolor{gsdnotebar}{HTML}{495057}

% ---------- Callout boxes ----------
\tcbset{
  gsdcalloutbase/.style={
    enhanced,
    breakable,
    left=10pt,
    right=8pt,
    top=6pt,
    bottom=6pt,
    boxrule=0pt,
    frame hidden,
    borderline west={4pt}{0pt}{#1},
    colback=#1,
    fonttitle=\bfseries\headingfont,
    fontupper=\small,
    arc=2pt
  }
}

\newtcolorbox{gsdtip}[1][Tip]{%
  gsdcalloutbase=gsdtipbg,
  borderline west={4pt}{0pt}{gsdtipbar},
  colback=gsdtipbg,
  title={#1}
}

\newtcolorbox{gsdwarning}[1][Warning]{%
  gsdcalloutbase=gsdwarnbg,
  borderline west={4pt}{0pt}{gsdwarnbar},
  colback=gsdwarnbg,
  title={#1}
}

\newtcolorbox{gsdimportant}[1][Important]{%
  gsdcalloutbase=gsddangerbg,
  borderline west={4pt}{0pt}{gsddangerbar},
  colback=gsddangerbg,
  title={#1}
}

\newtcolorbox{gsdnote}[1][Note]{%
  gsdcalloutbase=gsdnotebg,
  borderline west={4pt}{0pt}{gsdnotebar},
  colback=gsdnotebg,
  title={#1}
}

\newtcolorbox{gsdsuccess}[1][Success]{%
  gsdcalloutbase=gsdsuccessbg,
  borderline west={4pt}{0pt}{gsdsuccessbar},
  colback=gsdsuccessbg,
  title={#1}
}


% ---------- Hyperref (needs colors already defined) ----------
\usepackage[
  colorlinks=true,
  linkcolor=gsdnavy,
  urlcolor=gsdaccent,
  citecolor=gsdnavy,
  bookmarksnumbered=true
]{hyperref}

% ---------- Header / footer ----------
\pagestyle{fancy}
\fancyhf{}
\fancyhead[L]{\small\headingfont\@title}
\fancyhead[R]{\small\headingfont\leftmark}
\fancyfoot[L]{\small\headingfont GSD v1.35.0 — CC BY-SA 4.0}
\fancyfoot[R]{\small\headingfont Page \thepage\ of \pageref*{LastPage}}
\renewcommand{\headrulewidth}{0.4pt}
\renewcommand{\footrulewidth}{0.4pt}

\fancypagestyle{plain}{%
  \fancyhf{}
  \fancyfoot[L]{\small\headingfont GSD v1.35.0 — CC BY-SA 4.0}
  \fancyfoot[R]{\small\headingfont Page \thepage\ of \pageref*{LastPage}}
  \renewcommand{\headrulewidth}{0pt}
  \renewcommand{\footrulewidth}{0.4pt}
}

% ---------- Chapter / section styling ----------
\titleformat{\chapter}[display]
  {\headingfont\huge\bfseries\color{gsdnavy}}
  {\chaptertitlename\ \thechapter}{20pt}{\Huge}
\titleformat{\section}
  {\headingfont\Large\bfseries\color{gsdnavy}}
  {\thesection}{1em}{}
\titleformat{\subsection}
  {\headingfont\large\bfseries\color{gsdnavy}}
  {\thesubsection}{1em}{}

% ---------- Listings (code blocks) ----------
\lstdefinestyle{gsdcode}{
  basicstyle=\ttfamily\small,
  backgroundcolor=\color{gsdcodebg},
  frame=leftline,
  framerule=2pt,
  rulecolor=\color{gsdaccent},
  xleftmargin=1em,
  xrightmargin=0.5em,
  aboveskip=0.8em,
  belowskip=0.8em,
  breaklines=true,
  breakatwhitespace=true,
  showstringspaces=false,
  columns=flexible,
  keepspaces=true
}
\lstset{style=gsdcode}

% ---------- Semantic macros ----------
\newcommand{\cclicense}{%
  \textit{Released under CC BY-SA 4.0.\\
  Authored by Tibsfox. GSD itself is MIT-licensed — see gsd-build/get-shit-done.}%
}

\newcommand{\chapterheading}[1]{{\headingfont\LARGE\bfseries\color{gsdnavy}#1\par\medskip}}
\newcommand{\sectionheading}[1]{{\headingfont\Large\bfseries\color{gsdnavy}#1\par\smallskip}}
 in their
%  preamble (this file pulls in the color scheme transitively).
% ============================================================================

% ---------- Fonts ----------
\usepackage{fontspec}
\setmainfont{DejaVu Serif}[
  Ligatures=TeX,
  Scale=0.95,
  BoldFont={DejaVu Serif Bold},
  ItalicFont={DejaVu Serif Italic},
  BoldItalicFont={DejaVu Serif Bold Italic}
]
\newfontfamily\headingfont{DejaVu Sans}[Scale=0.95]
\newfontfamily\monofont{DejaVu Sans Mono}[Scale=0.88]
\setmonofont{DejaVu Sans Mono}[Scale=0.88]

% ---------- Page geometry ----------
\usepackage[
  letterpaper,
  top=1in,
  bottom=1in,
  left=1.25in,
  right=1.25in,
  headheight=14pt
]{geometry}

% ---------- Color + callout engine (must load BEFORE color-scheme input) ----
\usepackage{xcolor}
\usepackage[most]{tcolorbox}

% ---------- Typesetting packages ----------
\usepackage{tabularx}
\usepackage{longtable}
\usepackage{booktabs}
\usepackage{colortbl}
\usepackage{lastpage}
\usepackage{enumitem}
\usepackage{listings}
\usepackage{tikz}
\usepackage{fancyhdr}
\usepackage{titlesec}
\usepackage{parskip}
\usepackage{microtype}

% ---------- Pull in the color scheme (defines gsdnavy, gsdaccent, callouts) -
% ============================================================================
%  gsd-color-scheme.tex
%  Color palette + tcolorbox environment definitions for all GSD User Guides
% ============================================================================

\definecolor{gsdnavy}{HTML}{1B2A4A}
\definecolor{gsdaccent}{HTML}{CB3837}
\definecolor{gsdcodebg}{HTML}{F5F5F5}
\definecolor{gsdcodeborder}{HTML}{CCCCCC}

\definecolor{gsdtipbg}{HTML}{E8F4FD}
\definecolor{gsdtipbar}{HTML}{2B6CB0}

\definecolor{gsdwarnbg}{HTML}{FFF3CD}
\definecolor{gsdwarnbar}{HTML}{B78D00}

\definecolor{gsddangerbg}{HTML}{F8D7DA}
\definecolor{gsddangerbar}{HTML}{B02A37}

\definecolor{gsdsuccessbg}{HTML}{D4EDDA}
\definecolor{gsdsuccessbar}{HTML}{198754}

\definecolor{gsdnotebg}{HTML}{F0F0F5}
\definecolor{gsdnotebar}{HTML}{495057}

% ---------- Callout boxes ----------
\tcbset{
  gsdcalloutbase/.style={
    enhanced,
    breakable,
    left=10pt,
    right=8pt,
    top=6pt,
    bottom=6pt,
    boxrule=0pt,
    frame hidden,
    borderline west={4pt}{0pt}{#1},
    colback=#1,
    fonttitle=\bfseries\headingfont,
    fontupper=\small,
    arc=2pt
  }
}

\newtcolorbox{gsdtip}[1][Tip]{%
  gsdcalloutbase=gsdtipbg,
  borderline west={4pt}{0pt}{gsdtipbar},
  colback=gsdtipbg,
  title={#1}
}

\newtcolorbox{gsdwarning}[1][Warning]{%
  gsdcalloutbase=gsdwarnbg,
  borderline west={4pt}{0pt}{gsdwarnbar},
  colback=gsdwarnbg,
  title={#1}
}

\newtcolorbox{gsdimportant}[1][Important]{%
  gsdcalloutbase=gsddangerbg,
  borderline west={4pt}{0pt}{gsddangerbar},
  colback=gsddangerbg,
  title={#1}
}

\newtcolorbox{gsdnote}[1][Note]{%
  gsdcalloutbase=gsdnotebg,
  borderline west={4pt}{0pt}{gsdnotebar},
  colback=gsdnotebg,
  title={#1}
}

\newtcolorbox{gsdsuccess}[1][Success]{%
  gsdcalloutbase=gsdsuccessbg,
  borderline west={4pt}{0pt}{gsdsuccessbar},
  colback=gsdsuccessbg,
  title={#1}
}


% ---------- Hyperref (needs colors already defined) ----------
\usepackage[
  colorlinks=true,
  linkcolor=gsdnavy,
  urlcolor=gsdaccent,
  citecolor=gsdnavy,
  bookmarksnumbered=true
]{hyperref}

% ---------- Header / footer ----------
\pagestyle{fancy}
\fancyhf{}
\fancyhead[L]{\small\headingfont\@title}
\fancyhead[R]{\small\headingfont\leftmark}
\fancyfoot[L]{\small\headingfont GSD v1.35.0 — CC BY-SA 4.0}
\fancyfoot[R]{\small\headingfont Page \thepage\ of \pageref*{LastPage}}
\renewcommand{\headrulewidth}{0.4pt}
\renewcommand{\footrulewidth}{0.4pt}

\fancypagestyle{plain}{%
  \fancyhf{}
  \fancyfoot[L]{\small\headingfont GSD v1.35.0 — CC BY-SA 4.0}
  \fancyfoot[R]{\small\headingfont Page \thepage\ of \pageref*{LastPage}}
  \renewcommand{\headrulewidth}{0pt}
  \renewcommand{\footrulewidth}{0.4pt}
}

% ---------- Chapter / section styling ----------
\titleformat{\chapter}[display]
  {\headingfont\huge\bfseries\color{gsdnavy}}
  {\chaptertitlename\ \thechapter}{20pt}{\Huge}
\titleformat{\section}
  {\headingfont\Large\bfseries\color{gsdnavy}}
  {\thesection}{1em}{}
\titleformat{\subsection}
  {\headingfont\large\bfseries\color{gsdnavy}}
  {\thesubsection}{1em}{}

% ---------- Listings (code blocks) ----------
\lstdefinestyle{gsdcode}{
  basicstyle=\ttfamily\small,
  backgroundcolor=\color{gsdcodebg},
  frame=leftline,
  framerule=2pt,
  rulecolor=\color{gsdaccent},
  xleftmargin=1em,
  xrightmargin=0.5em,
  aboveskip=0.8em,
  belowskip=0.8em,
  breaklines=true,
  breakatwhitespace=true,
  showstringspaces=false,
  columns=flexible,
  keepspaces=true
}
\lstset{style=gsdcode}

% ---------- Semantic macros ----------
\newcommand{\cclicense}{%
  \textit{Released under CC BY-SA 4.0.\\
  Authored by Tibsfox. GSD itself is MIT-licensed — see gsd-build/get-shit-done.}%
}

\newcommand{\chapterheading}[1]{{\headingfont\LARGE\bfseries\color{gsdnavy}#1\par\medskip}}
\newcommand{\sectionheading}[1]{{\headingfont\Large\bfseries\color{gsdnavy}#1\par\smallskip}}

% ============================================================================
%  gsd-command-data.tex
%  Rendering macros for GSD command names, flags, files
%  Command inventory verified against github.com/gsd-build/get-shit-done
%  v1.35.0 (commit 66a5f93, released 2026-04-10)
% ============================================================================

% \gsdcmd{command-name}  — renders a GSD command name in accent color monospace
\newcommand{\gsdcmd}[1]{{\ttfamily\color{gsdaccent}/gsd-#1}}

% \gsdcolon{command-name} — renders with colon syntax (skill form)
\newcommand{\gsdcolon}[1]{{\ttfamily\color{gsdaccent}/gsd:#1}}

% \gsdflag{flag}  — renders a CLI flag in muted monospace
\newcommand{\gsdflag}[1]{{\ttfamily\footnotesize\color{gsdnotebar}-{}-#1}}

% \gsdfile{filename}  — renders a filename / path in italic monospace
\newcommand{\gsdfile}[1]{{\ttfamily\itshape #1}}

% \gsdkbd{key}  — renders a keyboard key
\newcommand{\gsdkbd}[1]{{\ttfamily\footnotesize\color{gsdnavy}[#1]}}

% ---------------------------------------------------------------------------
%  GSD v1.35.0 Command Inventory (73 commands, 16 categories)
%  Source: /tmp/gsd-upstream/docs/COMMANDS.md
%
%  Core Workflow:
%    /gsd-new-project        /gsd-new-workspace       /gsd-list-workspaces
%    /gsd-remove-workspace   /gsd-discuss-phase       /gsd-ui-phase
%    /gsd-plan-phase         /gsd-execute-phase       /gsd-verify-work
%    /gsd-next               /gsd-session-report      /gsd-ship
%    /gsd-ui-review          /gsd-audit-uat           /gsd-audit-milestone
%    /gsd-complete-milestone /gsd-milestone-summary   /gsd-new-milestone
%
%  Phase Management:
%    /gsd-add-phase          /gsd-insert-phase        /gsd-remove-phase
%    /gsd-list-phase-assumptions   /gsd-analyze-dependencies
%    /gsd-plan-milestone-gaps      /gsd-research-phase
%    /gsd-validate-phase
%
%  Navigation:
%    /gsd-progress  /gsd-resume-work  /gsd-pause-work  /gsd-manager  /gsd-help
%
%  Utility:
%    /gsd-explore    /gsd-undo       /gsd-import      /gsd-from-gsd2
%    /gsd-quick      /gsd-autonomous /gsd-do          /gsd-note
%    /gsd-debug      /gsd-add-todo   /gsd-check-todos /gsd-add-tests
%    /gsd-stats      /gsd-profile-user  /gsd-health   /gsd-cleanup
%
%  Diagnostics:
%    /gsd-forensics  /gsd-extract-learnings
%
%  Workstream Management:
%    /gsd-workstreams
%
%  Configuration:
%    /gsd-settings   /gsd-set-profile
%
%  Brownfield:
%    /gsd-map-codebase  /gsd-onboard  /gsd-migrate
%
%  AI Integration:
%    /gsd-ai-integration-phase  /gsd-eval-review
%
%  Update:
%    /gsd-update  /gsd-scan
%
%  Code Quality:
%    /gsd-code-review  /gsd-code-review-fix  /gsd-secure-phase
%
%  Fast & Inline:
%    /gsd-fast  /gsd-quick  /gsd-inline
%
%  Backlog & Thread:
%    /gsd-add-backlog  /gsd-review-backlog  /gsd-thread
%
%  State Management:
%    /gsd-pr-branch  /gsd-reapply-patches  /gsd-plant-seed  /gsd-trace
%
%  Community:
%    /gsd-join-discord
% ---------------------------------------------------------------------------


\begin{document}

% ============================================================
% Title page
% ============================================================
\begin{titlepage}
  \centering
  \vspace*{2cm}
  {\headingfont\Huge\bfseries\color{gsdnavy} GSD\par}
  \vspace{0.4cm}
  {\headingfont\LARGE The Beginner's Guide\par}
  \vspace{1.2cm}
  {\large\itshape Build real projects with AI ---\\
  no coding experience required.\par}
  \vspace{2cm}
  {\large For GSD v1.35.0\par}
  \vfill
  {\large Tibsfox\par}
  \vspace{1cm}
  \cclicense
  \vspace{1cm}
\end{titlepage}

\frontmatter
\pagenumbering{roman}
\tableofcontents
\cleardoublepage

\mainmatter
\pagenumbering{arabic}

% ============================================================
\chapter{Who This Guide Is For}
% ============================================================

Welcome. If you've ever had an idea for a small tool, a personal website, a
script that organizes your photos, or a little game---and you stopped because
you thought \emph{``I'd need to learn to code first''}---this guide is for you.

You don't need to learn to code. You need a good partner. GSD makes
Claude, Anthropic's AI assistant, into that partner.

GSD stands for \emph{Get Shit Done}. It's a small, free add-on for AI coding
assistants that keeps the AI organized while it builds what you asked for.
That's the whole pitch. You describe what you want in plain language. GSD
helps the AI think carefully, plan the work, do the work, and check its own
results. You stay in the driver's seat, but you don't have to be the mechanic.

\section{What Kind of Projects GSD Helps With}

GSD is happiest with software-shaped projects: websites, small apps,
browser extensions, scripts that automate a chore, a tool that takes a
spreadsheet and makes something useful out of it. It also does well with
\emph{creative} software: a little game, an interactive story, a digital
zine, an art generator.

It can help you fix a bug in something you already have. It can help you
read code someone else wrote and explain it to you. It can help you learn
by doing.

It is \emph{not} a magic wand that builds anything you can imagine with
zero effort. You still need to know what you want, and you still need to
answer questions the assistant asks you along the way. That part is good
news: the better you describe your idea, the better the result.

\section{What You Don't Need to Know}

\begin{itemize}
  \item You do not need to know how to program.
  \item You do not need to know what ``git'' is.
  \item You do not need to have used a terminal before.
  \item You do not need a computer science degree.
  \item You do not need to understand what's happening ``under the hood.''
\end{itemize}

You \emph{will} be typing a few commands into a window, and I'll walk you
through every one of them. I'll also tell you what each command does in plain
English before you run it.

\begin{gsdtip}[A gentle promise]
If something in this guide uses a word you don't know, I'll define it on the
spot. If a chapter feels confusing, it's the chapter's fault, not yours.
Skim ahead, come back, re-read. The commands are forgiving.
\end{gsdtip}

\section{What You'll Be Able to Do by the End}

By the time you reach the last page, you will have:

\begin{enumerate}
  \item Installed GSD on your computer.
  \item Started a real project with \gsdcmd{new-project}.
  \item Watched the assistant build and check its own work.
  \item Used \gsdcmd{progress} and \gsdcmd{next} to navigate the project.
  \item Fixed things that weren't quite right.
  \item Finished a milestone and started a new one.
\end{enumerate}

That's a full loop. Once you've done it once, you've done it. Everything
else is variations on the theme.

\bigskip
\noindent In the next chapter we'll set everything up. It takes about
fifteen minutes.

% ============================================================
\chapter{Setting Up}
% ============================================================

You need four things. Three of them you probably already have.

\section{What You'll Need}

\begin{enumerate}
  \item \textbf{A computer.} Mac, Windows, or Linux. GSD works on all three.
  \item \textbf{An internet connection.} The assistant talks to Anthropic's
        servers to think.
  \item \textbf{Claude Code.} This is Anthropic's coding assistant. It runs
        on your computer and is the thing GSD plugs into. You can get it at
        \texttt{claude.com/claude-code}. A free trial is usually available;
        a paid plan unlocks more.
  \item \textbf{Node.js.} This is a free tool that lets your computer run
        the GSD installer. You only install it once.
\end{enumerate}

\section{Installing Node.js}

Go to \texttt{nodejs.org}, click the big green ``LTS'' download button, run
the installer, click through the defaults. That's it.

\begin{gsdnote}[What \emph{is} Node.js, really?]
Node.js is a piece of software that runs small programs written in a
language called JavaScript. GSD's installer is one of those small programs.
You will not be writing any JavaScript yourself. Node.js is just the engine
that starts the installer.
\end{gsdnote}

To check that it worked, open a terminal window.

\begin{itemize}
  \item \textbf{Mac:} press \gsdkbd{Cmd}+\gsdkbd{Space}, type
        \texttt{terminal}, press \gsdkbd{Enter}.
  \item \textbf{Windows:} open the Start menu, type
        \texttt{powershell}, press \gsdkbd{Enter}.
  \item \textbf{Linux:} you already know.
\end{itemize}

A black (or light, if that's your taste) window opens with a blinking cursor.
Type:

\begin{lstlisting}
node --version
\end{lstlisting}

Press \gsdkbd{Enter}. You should see something like \texttt{v20.11.0}. If
you do, Node.js is installed and you can close the terminal.

\begin{gsdtip}[Version numbers]
The exact numbers don't matter for this guide, as long as the first number
is at least 18. If you see \texttt{v18}, \texttt{v20}, or \texttt{v22},
you're fine.
\end{gsdtip}

\section{Running the GSD Installer}

With Node.js installed, installing GSD is a single command. Open Claude
Code (it gives you its own terminal prompt) and type:

\begin{lstlisting}
npx get-shit-done-cc@latest
\end{lstlisting}

Press \gsdkbd{Enter}. The installer will ask you two questions.

\section{What the Installer Asks}

\textbf{Question 1: Which runtime?} The installer lists several AI coding
assistants it supports. Pick \textbf{Claude Code}. (If you're using a
different assistant, pick that one---the rest of this guide still applies;
only the way you type commands changes very slightly.)

\textbf{Question 2: Global or local?} Pick \textbf{Global}. ``Global''
means GSD is available in every project on your computer. ``Local'' means
only in the current folder. For a first project, global is easier.

The installer will print a few lines about copying files, then finish.
That's it. GSD is installed.

\section{Verifying It Worked}

In Claude Code, type:

\begin{lstlisting}
/gsd-help
\end{lstlisting}

Press \gsdkbd{Enter}. You should see a list of GSD commands scroll past.
If you do, congratulations---GSD is working. If you don't, see
Appendix~A for troubleshooting.

\begin{gsdsuccess}[You're installed]
The hardest part of this whole guide is now behind you. Everything from
here on is just talking to your assistant.
\end{gsdsuccess}

\bigskip
\noindent Next up: a short chapter on \emph{what} GSD is actually doing,
so the commands in Chapter~4 feel like friends instead of strangers.

% ============================================================
\chapter{What Is GSD?}
% ============================================================

You don't strictly need to read this chapter. You could skip to Chapter~4,
run the commands, and have a project. But if you understand \emph{why} GSD
exists, the commands will feel obvious instead of arbitrary. It's a short
chapter. Stick with me.

\section{The Problem: AI Forgets Mid-Conversation}

If you've used an AI chatbot for anything longer than a quick question,
you've probably noticed something strange: it gets worse as the
conversation gets longer. It forgets what you said earlier. It contradicts
itself. It starts giving vaguer answers. People who use AI assistants for
real work have a nickname for this: \emph{``context rot.''}

It happens because an AI assistant has a limited ``conversation memory''---
the amount of back-and-forth it can hold in its head at once. When the
conversation fills up that memory, old things start getting squeezed out
or blurred. Try to build a real project in one long chat and the
assistant will lose the thread.

\section{The Solution: GSD Keeps the AI Organized}

GSD solves the memory problem in a clever way: instead of keeping one
long conversation, it writes everything important down into files, and
starts fresh conversations when it needs to. Each fresh conversation is
given exactly the notes it needs and nothing more.

Think of it like this: imagine you hired a very skilled contractor to
build a treehouse, but the contractor had a medical condition that made
him forget everything at lunchtime. The only way to get the treehouse
built is to:

\begin{enumerate}
  \item Write down the plan in the morning while he still remembers the
        conversation.
  \item After lunch, hand him the written plan and say ``build exactly
        this.''
  \item Check what he built against the plan at the end of the day.
\end{enumerate}

GSD is the clipboard. The assistant is the contractor. You are the
homeowner who hired him. GSD keeps notes, hands out plans, and checks work,
so the assistant can stay focused on one small task at a time and do it well.

\section{The Workflow, As a Story}

Here's the whole GSD workflow in one paragraph:

\begin{quote}
\emph{You describe what you want. GSD asks you questions until it
understands. It writes down everything it learned---what you want, how
it's going to build it, in what order. It breaks the work into phases.
For each phase: it makes a plan, it does the work, it checks the work.
When all phases are done, you ship a version. When you want to add
more, you start a new version.}
\end{quote}

That's it. That's GSD. The commands in Chapter~4 each correspond to
one step of that story.

\section{What ``Spec-Driven Development'' Means}

You might see the phrase \emph{spec-driven development} in GSD's
documentation. It sounds jargony but the idea is homey: before you build
something, write down what you're building. The written description is
called a ``spec'' (short for specification). Spec-driven development just
means: write the spec first, then build from the spec.

GSD writes the spec \emph{with you}, by asking questions. You don't have
to know how to write a spec. You just have to answer the questions
honestly.

\begin{gsdnote}[The spec lives in your project]
When GSD writes a spec, it lives in a folder called \gsdfile{.planning/}
inside your project. You don't normally open those files yourself, but
they're yours. They're yours to read, yours to keep, yours to hand to
another developer or another AI later.
\end{gsdnote}

\bigskip
\noindent Enough theory. Let's build something.

% ============================================================
\chapter{Your First Project --- A Complete Walkthrough}
% ============================================================

This is the longest chapter in the guide. Take your time. If you follow
along on your own computer, you'll have a real project by the end.

For the walkthrough I'll pretend we're building a small personal website:
a one-page site with a bio, a list of projects, and a contact link. Pick
your own idea if you prefer---the steps are identical.

\section{Start a Project Folder}

Before you run any GSD commands, make a folder for your project. In the
terminal:

\begin{lstlisting}
mkdir my-website
cd my-website
\end{lstlisting}

\texttt{mkdir} means ``make directory'' (a directory is the same thing as
a folder). \texttt{cd} means ``change directory''---it's how you walk
into the folder you just made. After \texttt{cd my-website}, your
terminal is ``inside'' that folder.

\begin{gsdtip}[Pick a sensible name]
Use lowercase letters and dashes instead of spaces: \texttt{my-website},
\texttt{photo-tagger}, \texttt{birthday-quiz}. It makes everything easier
later.
\end{gsdtip}

Now open Claude Code in this folder. (In most setups, Claude Code picks
up the folder you started it in.)

\section{The First Command: \gsdcmd{new-project}}

This command is how you tell GSD ``I have an idea and I want to start a
project.'' Type:

\begin{lstlisting}
/gsd-new-project
\end{lstlisting}

Press \gsdkbd{Enter}. GSD wakes up and starts asking questions.

\section{What GSD Asks You, and Why}

GSD's first job is to understand your idea well enough to build it. It
does this by asking a series of questions. The exact questions vary, but
they usually cover things like:

\begin{itemize}
  \item What is the project? (One or two sentences is fine.)
  \item Who is it for? (Yourself? Your friends? Strangers on the internet?)
  \item What should it do when you're done? (Three or four bullet points.)
  \item Are there things that are explicitly \emph{out of scope}? (Things
        you do \emph{not} want, so the assistant doesn't gold-plate.)
  \item Is there anything it should look like or feel like? (A mood, a
        color, another site you admire.)
\end{itemize}

Answer each question the way you'd answer a friend. If you're not sure
about an answer, say ``I'm not sure, what do you recommend?''---the
assistant will offer options.

\begin{gsdtip}[This is the most important step]
The better you answer these questions, the better the final project.
People who rush this step end up with projects that miss the mark. People
who take their time end up surprised by how good the result is. Budget
ten or fifteen minutes for the first question-and-answer session. You
won't regret it.
\end{gsdtip}

\section{What Gets Created}

When the question session finishes, GSD will write a handful of files
into a folder called \gsdfile{.planning/} inside your project. The
important ones are:

\begin{itemize}
  \item \gsdfile{.planning/PROJECT.md} --- the summary of your idea. The
        thing you just described, written up cleanly.
  \item \gsdfile{.planning/REQUIREMENTS.md} --- a more detailed breakdown
        of what ``done'' looks like.
  \item \gsdfile{.planning/ROADMAP.md} --- the plan. A list of phases
        (numbered 1, 2, 3\ldots) that together add up to your project.
  \item \gsdfile{.planning/STATE.md} --- a living status file. GSD updates
        it as you progress.
\end{itemize}

You don't have to read any of these files. But if you're curious, open
\gsdfile{PROJECT.md} in any text editor and read it. It should feel like
your idea, written back to you.

\begin{gsdimportant}[Don't edit these files by hand]
If you want to change something in the plan---say you forgot to mention
a feature---don't open \gsdfile{ROADMAP.md} and type into it. Instead,
tell the assistant: \emph{``I forgot to mention I want a dark mode
toggle. Can you update the roadmap?''} It will update the files
correctly and keep everything consistent.
\end{gsdimportant}

\section{Understanding the Roadmap}

The roadmap is the most interesting of the files. It divides your
project into \textbf{phases}. A phase is a bite-sized chunk of work that
ends in something working. For a personal website, GSD might generate
phases like:

\begin{enumerate}
  \item \textbf{Phase 1:} Project skeleton and layout.
  \item \textbf{Phase 2:} Bio and about section.
  \item \textbf{Phase 3:} Projects list.
  \item \textbf{Phase 4:} Contact section and polish.
\end{enumerate}

Each phase is completed on its own. You finish Phase 1 (and \emph{see}
it working) before Phase 2 starts. This is how GSD keeps itself honest:
small chunks, visible progress, no six-week silent marathon that ends
in a pile of broken files.

\begin{gsdnote}[Phases are not sacred]
If you start the first phase and realize the plan is wrong, you can
stop and ask the assistant to revise the roadmap. Plans are meant to
change.
\end{gsdnote}

\section{Checking Where You Are}

At any point, you can ask GSD to tell you where you are in the project:

\begin{lstlisting}
/gsd-progress
\end{lstlisting}

This prints a small status report: which phase you're on, what's done,
what's next. It's a friendly ``you are here'' sign, and you'll use it
constantly.

\bigskip
\noindent You now have a project, a plan, and a roadmap. Nothing has
been \emph{built} yet---that's what the next chapter is for.

% ============================================================
\chapter{Building Your First Phase}
% ============================================================

GSD builds your project one phase at a time. In this chapter we'll walk
through the three commands that make up a single phase: plan, execute,
verify.

\section{What a ``Phase'' Is, Again}

A phase is a chunk of work with a clear beginning and end. For your
first phase, GSD will pick something small enough to finish quickly but
big enough that you can see real progress. In our website example,
Phase 1 is usually just ``set up the skeleton'': a page that loads and
shows ``Hello, world'' or similar.

Each phase goes through three steps:

\begin{enumerate}
  \item \textbf{Plan} --- GSD researches and writes a detailed plan for
        the phase.
  \item \textbf{Execute} --- GSD does the work in the plan.
  \item \textbf{Verify} --- GSD checks the work and asks you to
        double-check.
\end{enumerate}

Let's go.

\section{Step One: \gsdcmd{plan-phase}}

Type:

\begin{lstlisting}
/gsd-plan-phase 1
\end{lstlisting}

Press \gsdkbd{Enter}. (The \texttt{1} at the end means ``phase 1.'' If
you leave it off, GSD uses the next phase that needs planning, which is
also fine for a first-time run.)

GSD now goes off and thinks. It may take a few minutes. It might ask you
a clarifying question or two. When it's done, it will have written a
detailed plan file for the phase into the \gsdfile{.planning/} folder.
You don't have to read it, but if you do, you'll see something
refreshingly clear: here's what we'll build, here's what each part does,
here's how we'll know it worked.

\begin{gsdtip}[Feel free to review the plan]
Before running the next command, you can always ask the assistant:
\emph{``Can you summarize the plan for me in plain English?''} It will
give you a short, readable summary. If anything sounds wrong, say so
before moving on.
\end{gsdtip}

\section{Step Two: \gsdcmd{execute-phase}}

Once the plan exists, ask GSD to build it:

\begin{lstlisting}
/gsd-execute-phase 1
\end{lstlisting}

This is the fun part. The assistant will start creating files, writing
code, and generally making things happen. You'll see a running log of
what it's doing. At some points it will stop and ask for your permission
before doing something potentially important---for example, before
running a command that changes files in a significant way.

\begin{gsdimportant}[About permission prompts]
Claude Code will occasionally pause and ask you to approve an action.
Always read the prompt before approving. The default behavior---where
the assistant asks first---is the safe setting, and it is the correct
setting for beginners. You should leave it alone. You'll see suggestions
online telling you to turn the prompts off for convenience. Don't. Those
prompts are your seat belt.
\end{gsdimportant}

\begin{gsdwarning}[A note on ``skip permissions'' flags]
You may come across a command-line flag in other tutorials that
completely disables permission prompts. Do not use it. It exists for
experienced developers running automated pipelines in throwaway
environments, and it bypasses every safety check the assistant has. For
everything in this guide, the default interactive mode is what you
want.
\end{gsdwarning}

When the phase finishes, GSD will usually have created several new
files in your project. Congratulations---something real exists now.

\section{Step Three: \gsdcmd{verify-work}}

Before moving on, ask GSD to check its own work:

\begin{lstlisting}
/gsd-verify-work 1
\end{lstlisting}

This runs a short self-review. GSD looks at what was built, compares
it to what the plan said should happen, runs any tests that were
written, and reports back. If everything passed, you'll see a cheerful
summary. If something is off, GSD will tell you exactly what and
suggest a fix plan.

\begin{gsdsuccess}[Phase 1 complete]
When verify-work reports success, the phase is done. Seriously---pause
for a moment to notice that a real piece of software now exists on your
computer that didn't exist an hour ago. You made that happen. It's a
small thing, and also it isn't.
\end{gsdsuccess}

\section{What to Do If Something Isn't Right}

Sometimes the phase finishes and the result doesn't match what you
had in your head. That's normal. Here's the order of things to try:

\begin{enumerate}
  \item \textbf{Describe what's wrong.} Tell the assistant, in plain
        English, what you expected versus what you got. ``The page
        loads, but the heading says `Hello world' and I wanted it to
        say my name.''
  \item \textbf{Let it fix.} The assistant will usually propose a
        small fix and, with your permission, apply it.
  \item \textbf{Re-run verify.} Type \gsdcmd{verify-work} again to
        confirm the fix worked.
  \item \textbf{If you're stuck,} try \gsdcmd{debug} followed by a
        description: \texttt{/gsd-debug "Heading doesn't update when
        I change my name"}. GSD has a systematic debugging mode that
        walks through the problem step by step.
\end{enumerate}

If you ever feel truly lost, Appendix~A has more help.

\bigskip
\noindent You now know the full rhythm: plan, execute, verify, repeat
for the next phase. The next chapter covers the little commands that
help you find your way around a project in progress.

% ============================================================
\chapter{Navigating Your Project}
% ============================================================

A GSD project grows as you work on it. After a few phases, you'll have
a handful of plan files, some finished code, maybe a note or two. This
chapter is about the short, reassuring commands that tell you where
you are and what to do next.

\section{\texorpdfstring{\gsdcmd{progress}}{/gsd-progress} --- ``Where am I?''}

Run this any time you want a bird's-eye view:

\begin{lstlisting}
/gsd-progress
\end{lstlisting}

You'll see a summary of the project, which phase is current, what's
complete, and what's on deck. This is the command to run when you
come back to the project after a break and have no idea what you
were doing.

\section{\texorpdfstring{\gsdcmd{next}}{/gsd-next} --- ``What should I do next?''}

If \gsdcmd{progress} says ``where am I,'' \gsdcmd{next} says ``what
next.'' Run it and GSD figures out the correct next step for your
project and either does it or tells you the command to run:

\begin{lstlisting}
/gsd-next
\end{lstlisting}

If Phase 1 is planned but not executed, \gsdcmd{next} will execute it.
If Phase 2 hasn't been planned, it will plan it. If every phase is
done, it will suggest completing the milestone. It's the ``just keep
going'' button.

\section{\texorpdfstring{\gsdcmd{help}}{/gsd-help} --- ``What commands exist?''}

If you forget which commands GSD offers:

\begin{lstlisting}
/gsd-help
\end{lstlisting}

You'll see the full list. You won't need most of them for a while,
but it's nice to know they're there.

\section{The \gsdfile{.planning/} Directory}

I mentioned this folder earlier; it's worth a section of its own
because you'll see it a lot. \gsdfile{.planning/} is where GSD keeps
the project's brain: the spec, the roadmap, the plans for each phase,
the status file, notes. It lives inside your project folder, right
alongside the code GSD writes.

\begin{gsdimportant}[Do not delete this folder]
If you delete \gsdfile{.planning/}, GSD forgets everything about your
project. Not the code---the code is safe---but everything about
\emph{why} the code is the way it is, what's planned next, and what's
been decided. The folder is small (usually a few hundred kilobytes at
most), so there's no reason to delete it.
\end{gsdimportant}

\begin{gsdnote}[If your project is in version control]
If you're using version control (Git) with your project, the
\gsdfile{.planning/} folder can be committed along with the code. It
means the next person to open the project (including future-you on a
different computer) gets the plan, not just the code.
\end{gsdnote}

\section{Pausing and Resuming}

If you have to stop mid-work (lunch, the end of the day, the cat
pressed a key), GSD has two commands to make coming back painless:

\begin{lstlisting}
/gsd-pause-work
/gsd-resume-work
\end{lstlisting}

\gsdcmd{pause-work} writes a short handoff note summarizing where
you are. \gsdcmd{resume-work} reads it back and brings the assistant
up to speed in a fresh session. Use them whenever you close your
laptop.

\bigskip
\noindent That's navigation. The next chapter covers the ``quick task''
mode for when you don't need the full phase workflow.

% ============================================================
\chapter{Quick Tasks}
% ============================================================

Not every task deserves a full roadmap. Sometimes you just want to
add a footer to your website, or rename a button, or tweak a color.
For those, GSD has a lighter command: \gsdcmd{quick}.

\section{When to Use \texorpdfstring{\gsdcmd{quick}}{/gsd-quick}}

Use it for small, self-contained tasks that you could describe in
one or two sentences. Things like:

\begin{itemize}
  \item ``Add a footer with my copyright to the bottom of every page.''
  \item ``Change the background color to a warm beige.''
  \item ``Fix this bug where the submit button does nothing.''
  \item ``Add a favicon.''
\end{itemize}

For anything that feels bigger than a sentence or two---a new
section, a new feature, a rethink of the layout---use the full
planning workflow (Chapter~4 and 5) instead.

\section{How to Use It}

\begin{lstlisting}
/gsd-quick
\end{lstlisting}

GSD will ask you what you want done. Describe it in plain English.
It will propose a short plan, ask your permission to proceed, and
then do the work.

\begin{gsdtip}[Quick doesn't mean sloppy]
Even in quick mode, GSD still plans first and verifies afterward---
just with less ceremony. Your work is still safe.
\end{gsdtip}

If you'd like GSD to be a little more thorough for a particular
quick task, you can add flags:

\begin{itemize}
  \item \gsdflag{discuss} --- have a short pre-planning discussion
        first.
  \item \gsdflag{research} --- let GSD look things up before planning.
  \item \gsdflag{full} --- run the full plan-check and verification
        steps.
\end{itemize}

For example: \texttt{/gsd-quick --full}.

\bigskip
\noindent The next chapter is a collection of small habits that make
the whole experience smoother.

% ============================================================
\chapter{Tips for Success}
% ============================================================

Everyone who uses GSD well does the same small number of things.
They are not secrets and none of them are clever---they're just
habits. Here they are.

\section{Be Specific About What You Want}

``Make me a website'' gets you an average website. ``Make me a
one-page site for my grandmother's 80th birthday, with her photo,
a short tribute, a photo gallery of the grandkids, and a guestbook
where people can leave messages'' gets you \emph{that} website. The
more specific you are, the closer the result is to what you pictured.

This is the \#1 habit. It outweighs every other piece of advice in
this chapter combined.

\section{Let GSD Ask Questions}

\begin{sloppypar}
When GSD asks a question---during \gsdcmd{new-project}, during
planning, during execution---answer it. Don't skip it, don't say
``whatever,'' don't type ``idk.'' Every question it asks is because
it noticed a place where the result could go in two different
directions and it wants you to pick. If you don't pick, it picks
for you, and you may not like its choice.
\end{sloppypar}

\begin{gsdtip}[If you truly don't know]
Say so: \emph{``I don't know, which do you recommend and why?''} GSD
will give you a recommendation with reasons, and you can just agree
if the reasoning sounds good.
\end{gsdtip}

\section{Don't Edit \gsdfile{.planning/} Files by Hand}

I said this earlier but it bears repeating. If you want to change
the plan, talk to the assistant. Editing \gsdfile{.planning/} files
directly is like rewriting the blueprints while the contractor is
still reading them: confusing, and things will fall out of sync.

\section{Capture Ideas With \texorpdfstring{\gsdcmd{add-todo}}{/gsd-add-todo}}

You'll get ideas while the assistant is working on something else.
``Oh, I should add a search box.'' ``It would be nice if the images
had captions.'' Don't try to hold those in your head. Run:

\begin{lstlisting}
/gsd-add-todo "Add search box to the bio page"
\end{lstlisting}

GSD will save it to a list. Later, when you're ready, you can review
the list:

\begin{lstlisting}
/gsd-check-todos
\end{lstlisting}

You'll see everything you captured and can pick one to work on.

\section{Take Breaks}

This is less about GSD and more about you. A fresh set of eyes
catches things tired eyes miss. If something is going in circles,
step away for twenty minutes. When you come back, run
\gsdcmd{resume-work} and you'll be picked up exactly where you
left off.

\section{Celebrate Small Wins}

The first time a phase completes successfully, stop and look at
what you made. Open the thing in a browser, run the script, see
the output. \emph{Tell someone.} Building real things is a big
deal, and the whole point of GSD is that you get to do it.

\bigskip
\noindent The final chapter covers what happens when you finish
your first project and want to do more.

% ============================================================
\chapter{What's Next?}
% ============================================================

You finished the last phase. \gsdcmd{verify-work} says everything
checks out. What now?

\section{Complete the Milestone}

A ``milestone'' in GSD is a completed version of your project. When
all phases of a milestone are done, run:

\begin{lstlisting}
/gsd-complete-milestone
\end{lstlisting}

GSD will archive the planning documents for this milestone, update
a historical log of everything you've shipped, and (if you're using
version control) tag a release. It's your graduation ceremony for
that version of the project.

\section{Start a New Milestone}

When you're ready to add more---new features, a bigger redesign,
whatever---start a new milestone:

\begin{lstlisting}
/gsd-new-milestone
\end{lstlisting}

GSD will ask you what the new milestone is about, generate a new
roadmap for it, and you're back in the plan-execute-verify loop for
this new version. The old milestone is still there, archived, so
you can always look back.

\begin{gsdtip}[Versions are free]
Don't be shy about calling something ``done'' and starting a new
milestone. A project with five small completed milestones is healthier
than a project with one never-finished milestone. Ship early, ship
often, celebrate each one.
\end{gsdtip}

\section{When You're Ready for More}

This guide has covered the core workflow. GSD has many more commands
for special situations: bringing GSD into an existing codebase,
running phases in the background, reviewing code, generating tests,
digging into bugs, and more. Run \gsdcmd{help} any time to see the
full list.

Everything you learned here still applies. The extra commands are
just variations on the same rhythm: describe, plan, do, check.

\section{Getting Help}

GSD has an active community on Discord. If you're stuck, have a
question, or just want to share what you built, the people there are
friendly and have probably run into the same thing you're running
into. You can join from inside GSD:

\begin{lstlisting}
/gsd-join-discord
\end{lstlisting}

\begin{gsdsuccess}[You did it]
You started with an idea and you finished something real. That is
the entire point. Everything else is just practice. The next project
will be easier, the one after that easier still, and before long
you'll be surprised at what you can make.
\end{gsdsuccess}

% ============================================================
\appendix
\chapter{Troubleshooting}
% ============================================================

A short list of things that sometimes go wrong and what to do about
them.

\section{``Command Not Found''}

\textbf{Symptom:} you type \gsdcmd{help} (or any GSD command) and the
assistant says it doesn't recognize it.

\textbf{Fix:} close Claude Code completely and open it again. GSD
commands get picked up when the assistant starts, so a fresh start
usually does it. If that doesn't work, re-run the installer
(\texttt{npx get-shit-done-cc@latest}) and pick the same options as
before.

\section{``Something Went Wrong With the Project''}

\textbf{Symptom:} you get errors mentioning missing files, or
\gsdcmd{progress} reports something confusing about the project
state.

\textbf{Fix:} run the health check:

\begin{lstlisting}
/gsd-health
\end{lstlisting}

It will tell you what's wrong. If it offers to repair, let it:

\begin{lstlisting}
/gsd-health --repair
\end{lstlisting}

This fixes most common issues automatically.

\section{``I'm Completely Lost''}

\textbf{Symptom:} you stepped away from a project for a week and you
have no idea what you were doing.

\textbf{Fix:} two commands, in order:

\begin{lstlisting}
/gsd-resume-work
/gsd-progress
\end{lstlisting}

\gsdcmd{resume-work} brings the assistant up to speed; \gsdcmd{progress}
reminds you. Between them you'll be re-oriented in under a minute.

\section{``The Phase Produced Something I Don't Want''}

\textbf{Symptom:} \gsdcmd{execute-phase} finished, but the result is
wrong in some way.

\textbf{Fix:} first, describe the problem to the assistant in your
own words. Most issues resolve with a short conversation. If you'd
rather undo the phase entirely and start over:

\begin{lstlisting}
/gsd-undo --phase 1
\end{lstlisting}

This rolls back the phase's changes safely. You can then re-plan and
re-execute it.

\section{``Everything Is Fine, I'm Just Anxious''}

\textbf{Symptom:} things are working but you're worried you're doing
it wrong.

\textbf{Fix:} you're not. Run \gsdcmd{progress}, see that you're in
a valid state, take a breath, and keep going. The workflow is
forgiving. Mistakes are recoverable. You are allowed to be a
beginner.

\bigskip
\begin{center}
\textit{--- End of guide ---}
\end{center}

\end{document}
